\section*{Part V: Campus Master Planning and Architecture}
\addcontentsline{toc}{section}{Part V: Campus Master Planning and Architecture}


\subsection{35. Begin with an academic space program}

List the activities to be housed, users, capacities, hours, required adjacencies, services, safety, acoustics, flexibility and accessibility. A 40-person seminar room is not a wet lab or a compute room. Dormitories, machine shops, faculty offices, libraries, kitchens, sports and medical services have different costs and approval conditions. A master plan is an operational program translated into buildings, not a rendering translated into financing.

Ask faculty and operations to define a typical day and a peak day before designing floor area.


\subsection{36. Net usable area versus gross built area}

A facility's gross area includes circulation, stairs, walls, toilets, service shafts and plant rooms. Depending on design, usable academic area may be materially smaller. Model this as an explicit efficiency ratio instead of assuming every constructed square foot is instructional capacity. Lab and accessibility requirements can reduce the ratio, and these may be essential, not inefficiencies.

Specify the denominator for every \textnormal{₹}/square-foot cost: land, built-up, carpet, chargeable, constructed or fit-out. Mismatched denominators can create major budget errors.


\subsection{37. The density versus experience trade-off}

A dense campus can improve interaction and reduce infrastructure cost per learner; excessive density harms safety, feedback, concentration and wellbeing. A sprawling campus can look prestigious while multiplying transport, landscaping, utilities and maintenance expense. Cluster activities that benefit from spontaneous interaction while isolating noisy, hazardous or high-security functions.

The correct design optimizes academic work and livability under capital constraints, not visual grandeur.


\subsection{38. Growth corridors and future phases}

A well-designed master plan preserves structural and utility capacity for later phases without prematurely constructing them. Reserve expansion corridors, redundant paths where necessary, land for required services, and logical circulation. Each phase should be useful with its own operational fire-safety and utilities plan.

Do not construct unneeded buildings simply because the architect's phase-2 rendering looks incomplete. Staged completion must avoid forcing constant disruption of teaching and research.


\subsection{39. Infrastructure is more than buildings}

Budget distribution transformers, backup power, networking, water storage, treatment plants, wastewater, drainage, access roads, parking, fire infrastructure, security, accessibility, solid waste, laboratories' specialized ventilation, hazardous-waste handling, and maintenance workshops. Infrastructure often consumes substantial cost without producing attractive photographs.

The founder needs a full \textbf{site infrastructure schedule} and service-capacity assumptions at each phase. Ask what fails if the grid goes down for 8 hours or water supply fails for a week.


\subsection{40. Climate and resilience design}

Heat, monsoon flooding, extreme rainfall, water scarcity and grid reliability affect the facility's lifetime economics. Site orientation, shade, passive cooling, drainage, water reuse, efficient systems and backup power can reduce future vulnerability. Validate with location-specific engineering models rather than labeling a campus ``sustainable'' because it has a solar panel.

A campus designed for a century should accommodate changes in technology, climate and student behavior through modular infrastructure and reserves for renewal.


\subsection{41. Research laboratories create specialist constraints}

A laboratory may need vibration control, clean rooms, regulated storage, access containment, emergency eyewash, specialized ventilation, gases, power conditioning, fire suppression or environmental monitoring. These requirements differ drastically by field. Never assume a standard classroom building can be retrofitted cheaply for any future research purpose.

Use a scientifically qualified facility brief, specialist design and required safety permissions before procurement. For many emerging fields, outsourcing or shared lab access is the rational 1st move.


\subsection{42. Residential facilities change the institution}

Housing can strengthen community, enable extended projects, and attract out-of-region faculty. It also creates 24-hour security, student welfare, food service, maintenance, health, conduct and safeguarding obligations. A residential campus is a continuous service business, not merely an academic building with beds.

Assess real housing demand and local rental alternatives; consider 3rd-party managed housing or phased residences. Protect academic finances from opaque housing subsidies or related-party property deals.
